\begin{frame}{수렴 경계: 발산하는 $\alpha$ 는 정확히 계산된다}
  선형회귀는 $J$ 가 이차함수라 \textbf{수렴·발산 경계값을 정확히
  계산}할 수 있다. 업데이트는 매 스텝 ``현재점과 최솟값 $\boldsymbol{w}^{*}$ 의
  오차''에 고정 행렬을 곱하는 선형 연산:
  \[
  \boldsymbol{w}^{(t+1)} - \boldsymbol{w}^{*} = \left(I - \frac{\alpha}{m} X^TX\right)
  (\boldsymbol{w}^{(t)} - \boldsymbol{w}^{*})
  \]
  (증명은 2.2절 --- 여기서는 결과를 받아들이고 숫자를 본다.)
  \vskip0.6em
  장난감 데이터에서 $\frac{1}{m}X^TX$ 의 고유값이 $\approx 0.120$ 과
  $\approx 5.547$ 이므로, 매 스텝 오차가 고유방향별로
  $(1-\alpha\cdot0.120)$ 배, $(1-\alpha\cdot5.547)$ 배.
  두 방향 모두 줄려면 $|1-\alpha\cdot5.547| < 1$, 즉
  \[
  0 < \alpha < \frac{2}{5.547} \approx 0.361 \;= \alpha^{*}
  \]
\end{frame}
